%=============================================================================
% Chapter 3 Solutions
%=============================================================================
\newpage
\section*{Chapter 3 Solutions}
\addcontentsline{toc}{section}{Chapter 3 Solutions}

%-----------------------------------------------------------------------------
\subsection*{3.1: Unsigned Hex Subtraction}
\addcontentsline{toc}{subsection}{3.1: Unsigned Hex Subtraction}
\hypertarget{sol-3.1}{}
\markright{3.1}

Compute $\texttt{5ED4} - \texttt{07A4}$ as unsigned 16-bit hex:

\begin{align*}
  \texttt{5ED4}_{16} &= 5 \times 16^3 + 14 \times 16^2 + 13 \times 16 + 4
                      = 24{,}276 \\
  \texttt{07A4}_{16} &= 0 \times 16^3 + 7 \times 16^2 + 10 \times 16 + 4
                      = 1{,}956
\end{align*}

Subtract column by column from right to left:
\begin{itemize}
  \item Ones: $4 - 4 = 0$
  \item Sixteens: $D - A = 13 - 10 = 3$
  \item $16^2$: $E - 7 = 14 - 7 = 7$
  \item $16^3$: $5 - 0 = 5$
\end{itemize}

Result:
\[ \texttt{5ED4} - \texttt{07A4} = \textbf{\texttt{5730}} \]

Verification:
\[ 24{,}276 - 1{,}956 = 22{,}320 = 5 \times 16^3 + 7 \times 16^2
+ 3 \times 16 + 0 = \texttt{5730}_{16} \checkmark \]

%-----------------------------------------------------------------------------
\subsection*{3.2: Signed Hex Subtraction (Sign-Magnitude)}
\addcontentsline{toc}{subsection}{3.2: Signed Hex Subtraction (Sign-Magnitude)}
\hypertarget{sol-3.2}{}
\markright{3.2}

In 16-bit sign-magnitude, the MSB (bit 15) is the sign bit. For both values:
\begin{itemize}
  \item $\texttt{5ED4} = \texttt{0101\,1110\,1101\,0100}_2$.
        Sign bit $= 0$ (positive), magnitude $= \texttt{5ED4}_{16} = 24{,}276$.
        So the value is $+24{,}276$.
  \item $\texttt{07A4} = \texttt{0000\,0111\,1010\,0100}_2$.
        Sign bit $= 0$ (positive), magnitude $= \texttt{07A4}_{16} = 1{,}956$.
        So the value is $+1{,}956$.
\end{itemize}

Subtraction:
\[ +24{,}276 - (+1{,}956) = +22{,}320 \]

The result is positive. In sign-magnitude: sign bit $= 0$, magnitude:
\[ 22{,}320 = \texttt{5730}_{16} \]

Result: $\textbf{\texttt{5730}}$ (same as unsigned since both operands are
positive and the result is positive).

%-----------------------------------------------------------------------------
\subsection*{3.3: Hex to Binary Conversion}
\addcontentsline{toc}{subsection}{3.3: Hex to Binary Conversion}
\hypertarget{sol-3.3}{}
\markright{3.3}

Each hex digit maps to exactly 4 binary bits:

\begin{center}
\begin{tabular}{cccc}
\toprule
\texttt{5} & \texttt{E} & \texttt{D} & \texttt{4} \\
\midrule
\texttt{0101} & \texttt{1110} & \texttt{1101} & \texttt{0100} \\
\bottomrule
\end{tabular}
\end{center}

\[ \texttt{5ED4}_{16} = \texttt{0101\,1110\,1101\,0100}_2 \]

\medskip\noindent\textbf{Why hex is attractive for computers:}

Hexadecimal is attractive because each hex digit corresponds to exactly 4 binary
bits (since $16 = 2^4$). This makes conversion between hex and binary trivial --
no arithmetic is required, just a direct substitution of each digit. It provides
a compact, human-readable representation: a 32-bit value requires only 8 hex
digits instead of 32 binary digits. This is far more practical than decimal, where
the conversion to/from binary requires repeated division.

%-----------------------------------------------------------------------------
\subsection*{3.4: Unsigned Octal Subtraction}
\addcontentsline{toc}{subsection}{3.4: Unsigned Octal Subtraction}
\hypertarget{sol-3.4}{}
\markright{3.4}

Compute $4365_8 - 3412_8$ as unsigned 16-bit octal numbers.

Subtract column by column from right to left (base 8):
\begin{itemize}
  \item Ones: $5 - 2 = 3$
  \item Eights: $6 - 1 = 5$
  \item $8^2$: $3 - 4$. Since $3 < 4$, borrow 1 from the next column:
        $(3 + 8) - 4 = 7$, carry $= 1$.
  \item $8^3$: $4 - 3 - 1(\text{borrow}) = 0$
\end{itemize}

Result:
\[ 4365_8 - 3412_8 = \mathbf{0753_8} \]

Verification:
\begin{align*}
  4365_8 &= 4(512) + 3(64) + 6(8) + 5 = 2293 \\
  3412_8 &= 3(512) + 4(64) + 1(8) + 2 = 1802
\end{align*}
\[ 2293 - 1802 = 491 = 7(64) + 5(8) + 3 = 753_8 \checkmark \]

%-----------------------------------------------------------------------------
\subsection*{3.5: Signed Octal Subtraction (Sign-Magnitude)}
\addcontentsline{toc}{subsection}{3.5: Signed Octal Subtraction (Sign-Magnitude)}
\hypertarget{sol-3.5}{}
\markright{3.5}

In signed 12-bit sign-magnitude, the MSB (bit 11) is the sign bit.

Convert each octal number to binary (each octal digit is 3 bits):
\begin{itemize}
  \item $4365_8$: $4 = 100$, $3 = 011$, $6 = 110$, $5 = 101$.
        So $4365_8 = \texttt{100\,011\,110\,101}_2$ (12 bits).
        Sign bit (bit 11) $= 1$ (negative).
        Magnitude (bits 10--0) $= \texttt{00\,011\,110\,101}_2 = 245_{10}$.
        Value $= -245$.
  \item $3412_8$: $3 = 011$, $4 = 100$, $1 = 001$, $2 = 010$.
        So $3412_8 = \texttt{011\,100\,001\,010}_2$ (12 bits).
        Sign bit (bit 11) $= 0$ (positive).
        Magnitude (bits 10--0) $= \texttt{11\,100\,001\,010}_2 = 1794_{10}$.
        Value $= +1794$.
\end{itemize}

Subtraction:
\[ -245 - (+1794) = -2039 \]

Overflow check (Method 1): 12-bit sign-magnitude range is $[-(2^{11}-1),\; +(2^{11}-1)] =
[-2047,\; +2047]$. Since $|-2039| = 2039 \le 2047$, no overflow.

Encode $-2039$: sign bit $= 1$, magnitude $= 2039$.
\[ 2039_{10} = \texttt{11111110111}_2 \quad\text{(11 bits)} \]

Prepend the sign bit to form the full 12-bit representation:
\[ \texttt{1\,11111110111}_2 \]

Group into 3-bit chunks from the right and convert to octal:
\[ \underbrace{111}_{7}\;\underbrace{111}_{7}\;\underbrace{110}_{6}\;\underbrace{111}_{7} = 7767_8 \]

Result: $\mathbf{7767_8}$ (representing $-2039$ in sign-magnitude).

%-----------------------------------------------------------------------------
\subsection*{3.6: Unsigned 8-bit Subtraction Overflow}
\addcontentsline{toc}{subsection}{3.6: Unsigned 8-bit Subtraction Overflow}
\hypertarget{sol-3.6}{}
\markright{3.6}

\begin{flushright}
\begin{minipage}{0.92\linewidth}
\scriptsize
\textbf{Background: two ways to detect overflow}\par\smallskip
\begin{enumerate}
  \item \textbf{Decimal range check} --- compute the true result and verify
        it fits in the representable range.
  \item \textbf{Binary bit check} --- inspect carry/borrow or sign bits
        after the binary computation.
\end{enumerate}
\smallskip
Method~2 depends on the representation:

\smallskip
\begin{tabular}{@{}l l@{}}
\toprule
Representation & Binary overflow condition \\
\midrule
Unsigned & carry out (addition) or borrow out (subtraction) \\
Two's complement & sign flip: pos$+$pos$\to$neg or neg$+$neg$\to$pos \\
Sign-magnitude & magnitude exceeds available bits (no simple bit trick) \\
\bottomrule
\end{tabular}
\end{minipage}
\end{flushright}

\medskip
Unsigned 8-bit range: $[0, 255]$.

\[ 185 - 122 = 63 \]

Method 1: $63 \in [0, 255]$ $\checkmark$.
Method 2: $185 \ge 122$ (no borrow out) $\checkmark$.

\textbf{Neither overflow nor underflow}.

%-----------------------------------------------------------------------------
\subsection*{3.7: Signed Sign-Magnitude Subtraction Overflow}
\addcontentsline{toc}{subsection}{3.7: Signed Sign-Magnitude Subtraction Overflow}
\hypertarget{sol-3.7}{}
\markright{3.7}

In signed 8-bit sign-magnitude, bit 7 is the sign, bits 6--0 are the magnitude.
Range: $[-(2^7-1),\; +(2^7-1)] = [-127,\; +127]$.

Interpret the bit patterns:
\begin{itemize}
  \item $185 = 10111001_2$. Sign $= 1$ (negative), magnitude $= 0111001_2 = 57$.
        Value $= -57$.
  \item $122 = 01111010_2$. Sign $= 0$ (positive), magnitude $= 1111010_2 = 122$.
        Value $= +122$.
\end{itemize}

$185 - 122$ in sign-magnitude means:
\[ (-57) - (+122) = -179 \]

Method 1: $|-179| = 179 > 127$ (outside $[-127,\;+127]$) $\Rightarrow$
\textbf{overflow}.

(Method 2 does not apply to sign-magnitude---there is no simple bit trick
like two's complement sign flip. Just check the magnitude.)

\begin{flushright}
\begin{minipage}{0.92\linewidth}
\scriptsize
\textbf{Note: same bit pattern, three interpretations (3.6 vs.\ 3.7)}\par\smallskip
The bit pattern $185 = \texttt{10111001}_2$ has MSB $= 1$.
This is where the three representations diverge:

\smallskip
\begin{tabular}{@{}l l l@{}}
\toprule
Representation & 185 (\texttt{10111001}) & 122 (\texttt{01111010}) \\
\midrule
Unsigned & $185$ & $122$ \\
Sign-magnitude & $-57$ \;(sign=1, mag=$\texttt{0111001}=57$)
               & $+122$ \;(sign=0, mag=$\texttt{1111010}=122$) \\
Two's complement & $-71$ \;(flip+1: $\texttt{01000111}=71$)
                 & $+122$ \;(MSB=0, positive) \\
\bottomrule
\end{tabular}

\smallskip
Problems 3.1--3.2 used \texttt{5ED4} and \texttt{07A4} (both MSB $= 0$),
so all three representations gave the same value---the difference only
appears when MSB $= 1$.
\end{minipage}
\end{flushright}

%-----------------------------------------------------------------------------
\subsection*{3.8: Signed Sign-Magnitude Addition Overflow}
\addcontentsline{toc}{subsection}{3.8: Signed Sign-Magnitude Addition Overflow}
\hypertarget{sol-3.8}{}
\markright{3.8}

Using the same signed 8-bit sign-magnitude interpretations:
$185 = -57$, $122 = +122$.

$185 + 122$ in sign-magnitude means:
\[ (-57) + (+122) = +65 \]

Method 1: $65 \in [-127, +127]$ $\checkmark$ $\Rightarrow$ \textbf{no overflow}.

%-----------------------------------------------------------------------------
\subsection*{3.9: Two's Complement Saturating Addition}
\addcontentsline{toc}{subsection}{3.9: Two's Complement Saturating Addition}
\hypertarget{sol-3.9}{}
\markright{3.9}

\begin{flushright}
\begin{minipage}{0.92\linewidth}
\scriptsize
\textbf{Reading the problem statement}\par\smallskip
``signed 8-bit \textbf{decimal} integers stored in \textbf{two's complement} format''
--- each word has a distinct role:

\smallskip
\begin{tabular}{@{}l l@{}}
\toprule
Phrase & Meaning \\
\midrule
\textit{decimal} & The number is \textbf{written} in base-10 (notation) \\
\textit{8-bit} & The value is stored in 8 bits (storage size) \\
\textit{two's complement} & How to \textbf{interpret} the bit pattern (encoding) \\
\bottomrule
\end{tabular}

\smallskip
``decimal'' is \textbf{not} an encoding like octal or hex.
Octal means each digit $\to$ 3~bits, hex means each digit $\to$ 4~bits
--- these are notation-to-binary shortcuts.
``Decimal'' has no such shortcut; it simply means
``$151$ is written in base-10.''
The notation (decimal/octal/hex) and the storage size (8-bit/16-bit)
are independent concepts.

\smallskip
So the process is:
$151_{10} \to \texttt{10010111}_2$ (8-bit binary) $\to$
interpret as two's complement $\to -105$.
\end{minipage}
\end{flushright}

\begin{flushright}
\begin{minipage}{0.92\linewidth}
\scriptsize
\textbf{Two's complement: storage format vs.\ computation technique}\par\smallskip
Two's complement has two distinct roles that are easy to conflate:

\smallskip
\begin{tabular}{@{}l l@{}}
\toprule
Role & Description \\
\midrule
\textbf{Storage format} &
  MSB has weight $-2^{n-1}$.
  E.g.\ $\texttt{10010111} = -128+16+4+2+1 = -105$. \\
\textbf{Computation technique} &
  $A - B = A + (\sim\!B + 1)$. Subtraction via flip+1 and add. \\
\bottomrule
\end{tabular}

\smallskip
The computation technique works \textbf{regardless of the storage format}.
Example: IEEE 754 exponent $- 127$ (biased format, not two's complement):
\[
  127 = \texttt{01111111}
  \;\xrightarrow{\text{flip}}\;
  \texttt{10000000}
  \;\xrightarrow{+1}\;
  \texttt{10000001}
\]
Adding \texttt{10000001} to the stored exponent:
\[
  \text{stored} + \texttt{10000001}
  = \text{stored} + (2^8 - 127)
  = (\text{stored} - 127) + 2^8
\]
The $2^8$ part overflows out of 8~bits (carry out) and is discarded,
leaving $\text{stored} - 127$ --- the correct actual exponent.

\smallskip
\begin{tabular}{@{}l l l@{}}
\toprule
Context & Storage format & flip+1 used? \\
\midrule
MIPS signed integer & two's complement & yes (format \textit{and} technique) \\
Unsigned subtraction (3.1) & unsigned & yes (technique only) \\
FP exponent $-$ bias & biased & yes (technique only) \\
\bottomrule
\end{tabular}

\smallskip
Key point: flip+1 $\ne$ two's complement format.
Flip+1 is just a subtraction tool; any format can use it.
\end{minipage}
\end{flushright}

\medskip
Convert to 8-bit binary:
$151_{10} = \texttt{10010111}_2$, \;
$214_{10} = \texttt{11010110}_2$.
Both have MSB $= 1$, so both are negative in two's complement.

Add in binary:
\begin{verbatim}
    10010111   (151 unsigned, -105 in two's complement)
  + 11010110   (214 unsigned,  -42 in two's complement)
  ----------
  1 01101101
  ^ carry out (discard)
  = 01101101
\end{verbatim}

Overflow check (Method 2): neg $+$ neg $\to$ pos (MSB $= 0$)
$\Rightarrow$ \textbf{overflow}.

With \textbf{saturating arithmetic}, the result clamps to the most negative
representable value: $\mathbf{-128}$ ($= \texttt{10000000}_2$).

%-----------------------------------------------------------------------------
\subsection*{3.10: Two's Complement Saturating Subtraction}
\addcontentsline{toc}{subsection}{3.10: Two's Complement Saturating Subtraction}
\hypertarget{sol-3.10}{}
\markright{3.10}

$151_{10} = \texttt{10010111}_2$, \;
$214_{10} = \texttt{11010110}_2$.

Subtraction $= $ addition of flip+1:
\begin{verbatim}
  ~11010110 = 00101001
  00101001 + 1 = 00101010   (flip+1 of 214)
\end{verbatim}

Add:
\begin{verbatim}
    10010111   (151)
  + 00101010   (flip+1 of 214)
  ----------
    11000001
\end{verbatim}

Overflow check (Method 2): neg $+$ pos $\to$ \textbf{never overflows}.

Result: $\texttt{11000001}_2$. MSB $= 1$ (negative),
$= -128 + 64 + 1 = \mathbf{-63}$.
Saturating arithmetic has no effect.

%-----------------------------------------------------------------------------
\subsection*{3.11: Unsigned Saturating Addition}
\addcontentsline{toc}{subsection}{3.11: Unsigned Saturating Addition}
\hypertarget{sol-3.11}{}
\markright{3.11}

Unsigned 8-bit range: $[0, 255]$.

\[ 151 + 214 = 365 \]

Since $365 > 255$, overflow occurs. With saturating arithmetic, the result
clamps to the maximum: $\mathbf{255}$.
%-----------------------------------------------------------------------------
\subsection*{3.12: Multiplication Table (6-bit, Fig.~3.3)}
\addcontentsline{toc}{subsection}{3.12: Multiplication Table (6-bit, Fig.~3.3)}
\hypertarget{sol-3.12}{}
\markright{3.12}

Multiply unsigned 6-bit octal values: $62_8 \times 12_8$.

Convert:
\[
  62_8 = 110\,010_2 = 50_{10} \quad(\text{multiplicand}), \qquad
  12_8 = 001\,010_2 = 10_{10} \quad(\text{multiplier}).
\]
Expected product:
\[
  50 \times 10 = 500 = 764_8.
\]

\medskip\noindent\textbf{Algorithm (Fig.~3.3 -- first multiplier version):}
The multiplicand register is 12 bits wide (double width). The product register
is 12 bits, initialized to 0. The multiplier register is 6 bits. Each iteration:

\begin{enumerate}
  \item Test multiplier bit 0.
  \item If 1: add multiplicand to product. If 0: do nothing.
  \item Shift multiplicand left 1 bit.
  \item Shift multiplier right 1 bit.
  \item Repeat for 6 iterations (one per multiplier bit).
\end{enumerate}

\begin{center}
\begin{tabular}{clll}
\toprule
\textbf{Step} & \textbf{Multiplicand} & \textbf{Multiplier} & \textbf{Product} \\
\midrule
Init & \texttt{000000\,110010} & \texttt{001010} & \texttt{000000000000} \\
1 & (bit 0 = 0, no add) & & \\
  & \texttt{000001\,100100} & \texttt{000101} & \texttt{000000000000} \\
2 & (bit 0 = 1, add) & & \\
  & \texttt{000011\,001000} & \texttt{000010} & \texttt{000001100100} \\
3 & (bit 0 = 0, no add) & & \\
  & \texttt{000110\,010000} & \texttt{000001} & \texttt{000001100100} \\
4 & (bit 0 = 1, add) & & \\
  & \texttt{001100\,100000} & \texttt{000000} & \texttt{000111110100} \\
5 & (bit 0 = 0, no add) & & \\
  & \texttt{011001\,000000} & \texttt{000000} & \texttt{000111110100} \\
6 & (bit 0 = 0, no add) & & \\
  & \texttt{110010\,000000} & \texttt{000000} & \texttt{000111110100} \\
\bottomrule
\end{tabular}
\end{center}

Product:
\[
  \texttt{000111110100}_2 = 500_{10} = \mathbf{764_8} \checkmark
\]

%-----------------------------------------------------------------------------
\subsection*{3.13: Multiplication Table (8-bit, Fig.~3.5)}
\addcontentsline{toc}{subsection}{3.13: Multiplication Table (8-bit, Fig.~3.5)}
\hypertarget{sol-3.13}{}
\markright{3.13}

The problem says ``octal unsigned 8-bit integers 62 and 12,'' so we use the same
decimal values:
\[
  \text{multiplicand} = 62_8 = 50_{10}, \qquad
  \text{multiplier} = 12_8 = 10_{10}.
\]
As 8-bit values: multiplicand $= 00110010_2$, multiplier $= 00001010_2$.

\medskip\noindent\textbf{Algorithm (Fig.~3.5 -- optimized multiplier):}
The multiplicand register is 8 bits. The product register is 16 bits, with the
multiplier initially placed in the right half. Each iteration:
\begin{enumerate}
  \item Test product bit 0 (rightmost bit of multiplier portion).
  \item If 1: add multiplicand to left half of product.
  \item Shift entire product register right 1 bit.
  \item Repeat for 8 iterations.
\end{enumerate}

\begin{center}
\begin{tabular}{cll}
\toprule
\textbf{Step} & \textbf{Multiplicand} & \textbf{Product [HI $|$ LO/Multiplier]} \\
\midrule
Init & \texttt{00110010} & \texttt{00000000\,00001010} \\
1 & (bit 0 = 0, no add; shift right) & \texttt{00000000\,00000101} \\
2 & (bit 0 = 1, add \& shift) & add: \texttt{00110010\,00000101}; shift: \texttt{00011001\,00000010} \\
3 & (bit 0 = 0, no add; shift right) & \texttt{00001100\,10000001} \\
4 & (bit 0 = 1, add \& shift) & add: \texttt{00111110\,10000001}; shift: \texttt{00011111\,01000000} \\
5 & (bit 0 = 0, no add; shift right) & \texttt{00001111\,10100000} \\
6 & (bit 0 = 0, no add; shift right) & \texttt{00000111\,11010000} \\
7 & (bit 0 = 0, no add; shift right) & \texttt{00000011\,11101000} \\
8 & (bit 0 = 0, no add; shift right) & \texttt{00000001\,11110100} \\
\bottomrule
\end{tabular}
\end{center}

Product:
\[
  \texttt{0000000111110100}_2 = 500_{10} = \mathbf{764_8} \checkmark
\]

%-----------------------------------------------------------------------------
\subsection*{3.14: Multiply Time Calculation}
\addcontentsline{toc}{subsection}{3.14: Multiply Time Calculation}
\hypertarget{sol-3.14}{}
\markright{3.14}

8-bit integer, 4 time units per step. Each iteration has 3 steps: (1a) test and
conditionally add, (1b) shift multiplicand/product, (1c) shift multiplier.

For an 8-bit multiply, there are 8 iterations.

\medskip\noindent\textbf{Hardware (simultaneous shifts):}
Each iteration: 1 add + 1 shift step (two shifts done simultaneously) $= 2$ steps
$\times\, 4 = 8$ time units per iteration.
\[
  \text{Total} = 8 \times 8 = \mathbf{64} \text{ time units.}
\]

\medskip\noindent\textbf{Software (sequential shifts):}
Each iteration: 1 add + 2 shift steps (done sequentially) $= 3$ steps
$\times\, 4 = 12$ time units per iteration.
\[
  \text{Total} = 8 \times 12 = \mathbf{96} \text{ time units.}
\]

\medskip\noindent Both Fig.~3.3 and Fig.~3.5 have the same number of iterations
(equal to the bit width), so the times are the same for each approach.

%-----------------------------------------------------------------------------
\subsection*{3.15: Multiply Time with Stacked Adders}
\addcontentsline{toc}{subsection}{3.15: Multiply Time with Stacked Adders}
\hypertarget{sol-3.15}{}
\markright{3.15}

For an 8-bit multiply with a combinational array of adders: there are 7 rows of
adders (the first partial product needs no addition; partial products 2 through 8
each need one row of addition). However, the text says 31 adders for a 32-bit
multiplier, so for 8-bit this would be 7 adders stacked vertically.

Each adder takes 4 time units, and they are connected in series (the output of
one feeds the next).
\[
  \text{Total} = 7 \times 4 = \mathbf{28} \text{ time units.}
\]

This is significantly faster than the iterative approach because all partial
products are computed simultaneously in combinational logic.

%-----------------------------------------------------------------------------
\subsection*{3.16: Multiply Time (Fig.~3.7)}
\addcontentsline{toc}{subsection}{3.16: Multiply Time (Fig.~3.7)}
\hypertarget{sol-3.16}{}
\markright{3.16}

Fig.~3.7 shows the Wallace tree multiplier approach, which uses a tree of
carry-save adders to reduce partial products in parallel.

For an 8-bit multiplier, there are 8 partial products. A Wallace tree reduces
$n$ partial products using $\lceil \log_{1.5} n \rceil$ levels of carry-save
adders, plus one final carry-propagate adder.

For $n = 8$: the number of reduction levels is
\[
  \lceil \log_{1.5}(8) \rceil = \lceil 5.13 \rceil = 4
\]
carry-save adder levels (each taking roughly 2 time units for a carry-save add
since there is no carry propagation, but the problem states each adder takes 4
time units).

With each reduction level taking 4 time units, plus one final
carry-propagate add of 4 time units:
\[
  \text{Total} \approx (4 + 1) \times 4 = \mathbf{20} \text{ time units.}
\]

Alternatively, if using a balanced binary tree of standard adders:
\[
  \lceil \log_2 8 \rceil = 3 \text{ levels, each 4 time units} = 12 \text{ time units.}
\]
The exact answer depends on interpretation, but the key point is that the tree
structure achieves $O(\log n)$ depth vs.\ $O(n)$ for the sequential approach.

%-----------------------------------------------------------------------------
\subsection*{3.17: Shift-and-Add Multiplication}
\addcontentsline{toc}{subsection}{3.17: Shift-and-Add Multiplication}
\hypertarget{sol-3.17}{}
\markright{3.17}

Compute $\texttt{0x33} \times \texttt{0x55}$ using shifts and adds. Convert:
\[
  \texttt{0x33} = 51_{10} = 00110011_2, \qquad
  \texttt{0x55} = 85_{10} = 01010101_2.
\]

We decompose one operand into powers of 2. Using $51 = 2^6 - 2^4 + 2^2 - 2^0$
requires 3 shifts and 3 add/subtracts (6 operations).

A more efficient factoring uses $51 = 3 \times 17 = (2 + 1)(16 + 1)$:

\begin{verbatim}
    t1 = 0x55 + (0x55 << 4)   # 0x55 * 17  = 0x0595
    t2 = t1 + (t1 << 1)       # 0x55 * 51  = 0x10EF
\end{verbatim}

This uses only 2 shifts and 2 adds (4 operations total), which is optimal.

Verification:
\[
  51 \times 85 = 4335 = \texttt{0x10EF} \checkmark
\]

%-----------------------------------------------------------------------------
\subsection*{3.18: Division Table (Fig.~3.8)}
\addcontentsline{toc}{subsection}{3.18: Division Table (Fig.~3.8)}
\hypertarget{sol-3.18}{}
\markright{3.18}

Divide $74 \div 21$ using the first division algorithm (Fig.~3.8). Both are
unsigned 8-bit integers.
\[
  74 = 01001010_2, \qquad 21 = 00010101_2.
\]
Expected:
\[
  74 \div 21 = 3 \text{ remainder } 11 \qquad (74 = 21 \times 3 + 11).
\]

\medskip\noindent\textbf{Algorithm (Fig.~3.8):}
The divisor starts in the left half of a 16-bit register. The remainder register
is 16 bits, initialized with the dividend in the right half. The quotient is 8
bits, initialized to 0.

Each iteration (9 iterations for 8-bit, since we start with a subtract):
\begin{enumerate}
  \item Subtract divisor from remainder.
  \item If remainder $\ge 0$: shift quotient left, set new bit to 1.
  \item If remainder $< 0$: restore remainder (add divisor back), shift quotient
        left, set new bit to 0.
  \item Shift divisor right 1 bit.
\end{enumerate}

After 9 iterations, the quotient register holds the quotient and the left half
of the remainder register holds the remainder.

Final result:
\[
  \text{Quotient} = 3 = 00000011_2, \qquad \text{Remainder} = 11 = 00001011_2.
\]
\[
  \mathbf{74 \div 21 = 3} \text{ remainder } \mathbf{11}
\]

%-----------------------------------------------------------------------------
\subsection*{3.19: Division Table (Fig.~3.11)}
\addcontentsline{toc}{subsection}{3.19: Division Table (Fig.~3.11)}
\hypertarget{sol-3.19}{}
\markright{3.19}

Divide $74 \div 21$ using the optimized division hardware (Fig.~3.11).

In Fig.~3.11, the remainder and quotient share a single register (similar to how
Fig.~3.5 optimizes multiplication). The divisor is 8 bits. The
remainder/quotient register is 16 bits.

\medskip\noindent\textbf{Algorithm:}
Initialize the 16-bit remainder register with the dividend in the right half
(bits 7--0), left half zero. The divisor is in an 8-bit register.

\begin{enumerate}
  \item Shift the remainder register left 1 bit.
  \item Subtract the divisor from the left half of the remainder register.
  \item If the result $\ge 0$: set the rightmost bit (bit 0) of the remainder
        register to 1.
  \item If the result $< 0$: restore the left half (add divisor back), set bit 0
        to 0.
  \item Repeat for 8 iterations.
  \item After the loop, the left half contains the remainder (but shifted left by
        1 too many times). Shift the left half right by 1 to correct.
\end{enumerate}

After all steps: the right half holds the quotient, and the left half (after
correction) holds the remainder.

Final result:
\[
  \text{Quotient} = \mathbf{3}, \qquad \text{Remainder} = \mathbf{11}.
\]

Verification:
\[
  21 \times 3 + 11 = 63 + 11 = 74 \checkmark
\]

%-----------------------------------------------------------------------------
\subsection*{3.20: Bit Pattern as Integer}
\addcontentsline{toc}{subsection}{3.20: Bit Pattern as Integer}
\hypertarget{sol-3.20}{}
\markright{3.20}

$\texttt{0x0C000000}$ in binary:
\[
  \texttt{0000\,1100\,0000\,0000\,0000\,0000\,0000\,0000}_2
\]

\noindent\textbf{As unsigned integer:}
\[
  2^{27} + 2^{26} = 134{,}217{,}728 + 67{,}108{,}864 = \mathbf{201{,}326{,}592}
\]

\noindent\textbf{As two's complement signed integer:}
The MSB (bit 31) is 0, so the number is positive. The value is the same:
\[
  \mathbf{+201{,}326{,}592}
\]

%-----------------------------------------------------------------------------
\subsection*{3.21: Bit Pattern as MIPS Instruction}
\addcontentsline{toc}{subsection}{3.21: Bit Pattern as MIPS Instruction}
\hypertarget{sol-3.21}{}
\markright{3.21}

$\texttt{0x0C000000} = \texttt{0000\,1100\,0000\,0000\,0000\,0000\,0000\,0000}_2$.

Extract the opcode (bits 31--26):
\[
  \texttt{000011}_2 = 3_{10}
\]

Opcode 3 in MIPS is \texttt{jal} (jump and link), which uses J-type format:
\begin{verbatim}
  000011 00000000000000000000000000
  op=3   target address = 0
\end{verbatim}

The 26-bit target address field is all zeros. The full jump target is formed by:
\[
  \text{target} = \{\text{PC}[31{:}28],\; \texttt{00\,0000\,0000\,0000\,0000\,0000\,0000},\; 00\}
\]

The instruction is: \textbf{\texttt{jal 0x00000000}} (jump and link to address 0,
within the current 256 MB region defined by the upper 4 bits of PC+4).

%-----------------------------------------------------------------------------
\subsection*{3.22: Bit Pattern as Floating Point}
\addcontentsline{toc}{subsection}{3.22: Bit Pattern as Floating Point}
\hypertarget{sol-3.22}{}
\markright{3.22}

$\texttt{0x0C000000} = \texttt{0\;00011000\;00000000000000000000000}_2$.

IEEE 754 single precision fields:
\begin{itemize}
  \item Sign: $s = 0$ (positive).
  \item Exponent: $\texttt{00011000}_2 = 24_{10}$.
        Biased exponent with bias 127:
        \[ E = 24 - 127 = -103 \]
  \item Mantissa: $\texttt{00000000000000000000000}_2 = 0$.
        With hidden 1: significand $= 1.0$.
\end{itemize}

Value (since the exponent field is nonzero, this is a normal number):
\[
  (-1)^0 \times 1.0 \times 2^{-103} \approx \mathbf{9.8607613 \times 10^{-32}}
\]

This is a very small positive float.

%-----------------------------------------------------------------------------
\subsection*{3.23: IEEE 754 Single Precision Encoding}
\addcontentsline{toc}{subsection}{3.23: IEEE 754 Single Precision Encoding}
\hypertarget{sol-3.23}{}
\markright{3.23}

Convert $63.25$ to IEEE 754 single precision.

\noindent\textbf{Step 1:} Convert to binary.
\begin{align*}
  63 &= 111111_2 \\
  0.25 &= 0.01_2 \\
  63.25 &= 111111.01_2
\end{align*}

\noindent\textbf{Step 2:} Normalize.
\[
  111111.01_2 = 1.1111101_2 \times 2^5
\]

\noindent\textbf{Step 3:} Extract fields.
\begin{itemize}
  \item Sign: $s = 0$ (positive).
  \item Exponent: $E = 5$, biased $= 5 + 127 = 132 = 10000100_2$.
  \item Mantissa: $1111101\underbrace{0000000000000000}_{16\text{ zeros}}$
        (23 bits, drop the leading 1).
\end{itemize}

Binary representation:
\[
  \boxed{\texttt{0\;10000100\;11111010000000000000000}}
\]

Group into nibbles:
\[
  \texttt{0100\,0010\,0111\,1101\,0000\,0000\,0000\,0000} = \boxed{\texttt{0x427D0000}}
\]

%-----------------------------------------------------------------------------
\subsection*{3.24: IEEE 754 Double Precision Encoding}
\addcontentsline{toc}{subsection}{3.24: IEEE 754 Double Precision Encoding}
\hypertarget{sol-3.24}{}
\markright{3.24}

Same value $63.25 = 1.1111101_2 \times 2^5$.

\begin{itemize}
  \item Sign: $s = 0$.
  \item Exponent: $E = 5$, biased $= 5 + 1023 = 1028 = 10000000100_2$ (11 bits).
  \item Mantissa: $1111101$ followed by 45 zeros (52 bits total).
\end{itemize}

Binary representation:
\[
  \boxed{\texttt{0\;10000000100\;1111101}\underbrace{\texttt{000\ldots0}}_{45}}
\]

In hex:
\[
  \boxed{\texttt{0x404FA00000000000}}
\]

Grouping: $\texttt{0100\,0000\,0100\,1111\,1010\,0000\,\ldots\,0000}$.

%-----------------------------------------------------------------------------
\subsection*{3.25: IEEE 754 Negative Value Encoding}
\addcontentsline{toc}{subsection}{3.25: IEEE 754 Negative Value Encoding}
\hypertarget{sol-3.25}{}
\markright{3.25}

\[
  -1.5625 \times 10^{-1} = -0.15625
\]

\noindent\textbf{Step 1:} Convert to binary.
\[
  0.15625 = 0.00101_2 \quad\text{since } 2^{-3} + 2^{-5} = 0.125 + 0.03125 = 0.15625
\]

\noindent\textbf{Step 2:} Normalize.
\[
  0.00101_2 = 1.01_2 \times 2^{-3}
\]

\noindent\textbf{Step 3:} IEEE 754 single precision fields.
\begin{itemize}
  \item Sign: $s = 1$ (negative).
  \item Exponent: $E = -3$, biased $= -3 + 127 = 124 = 01111100_2$.
  \item Mantissa: $01\underbrace{000\ldots0}_{21}$ (23 bits).
\end{itemize}

Bit pattern:
\[
  \boxed{\texttt{1\;01111100\;01000000000000000000000}}
\]

Group into nibbles:
\[
  \texttt{1011\,1110\,0010\,0000\,\ldots\,0000} = \boxed{\texttt{0xBE200000}}
\]

%-----------------------------------------------------------------------------
\subsection*{3.26: DEC PDP-8 Floating Point Format}
\addcontentsline{toc}{subsection}{3.26: DEC PDP-8 Floating Point Format}
\hypertarget{sol-3.26}{}
\markright{3.26}

\[
  -0.15625 = -0.625 \times 2^{-2}
\]

The DEC PDP-8 format uses 36 bits total:
\begin{itemize}
  \item Bits 35--24: 12-bit exponent in two's complement.
  \item Bits 23--0: 24-bit fraction in two's complement. No hidden 1.
\end{itemize}

The number is represented as $\text{fraction} \times 2^{\text{exponent}}$, where the
fraction is a two's complement fixed-point value with the binary point immediately
after the sign bit. The fraction is normalized so that $0.5 \le |f| < 1$.

Normalizing: $-0.15625 = -\tfrac{5}{32} = -\tfrac{5}{8} \times 2^{-2} = -0.625 \times 2^{-2}$.

\noindent\textbf{Exponent:} $-2$ in 12-bit two's complement:
\[
  \texttt{111111111110}_2
\]

\noindent\textbf{Fraction:} $-0.625$ in 24-bit two's complement with the binary point
after the sign bit. The representation is $b_{23}.b_{22}b_{21}\ldots b_0$, where the
value is $-b_{23} + \sum_{i=0}^{22} b_i \cdot 2^{i-23}$.

For $-0.625$: the sign bit is 1, so we need $-1 + x = -0.625$, giving $x = 0.375 = 0.011_2$.
The 24 bits are therefore:
\[
  \texttt{101100000000000000000000}_2
\]

\medskip\noindent\textbf{36-bit pattern:}
\[
  \underbrace{\texttt{111111111110}}_{\text{exponent} = -2}\;
  \underbrace{\texttt{101100000000000000000000}}_{\text{fraction} = -0.625}
\]

\medskip\noindent\textbf{Comparison with IEEE 754:}
\begin{itemize}
  \item \textbf{Range:} The PDP-8 has a 12-bit two's complement exponent, giving
        a range of $2^{-2048}$ to $2^{+2047}$, far exceeding IEEE 754 single's
        $2^{-126}$ to $2^{+127}$.
  \item \textbf{Accuracy:} The PDP-8 has 24 fraction bits including the sign, so
        effectively 23 bits of magnitude --- comparable to IEEE 754 single (23-bit
        mantissa + hidden 1 = 24 bits of precision). The PDP-8 has slightly less
        precision since it lacks the hidden-1 optimization.
  \item The PDP-8 uses 36 bits total vs.\ IEEE 754 single's 32 bits.
\end{itemize}

%-----------------------------------------------------------------------------
\subsection*{3.27: IEEE 754-2008 Half Precision Format}
\addcontentsline{toc}{subsection}{3.27: IEEE 754-2008 Half Precision Format}
\hypertarget{sol-3.27}{}
\markright{3.27}

Format: 1 sign + 5-bit exponent (excess-16 bias) + 10-bit mantissa + hidden 1.
Total: 16 bits.

\[
  -0.15625 = -1.01_2 \times 2^{-3}
\]

\begin{itemize}
  \item Sign: $s = 1$.
  \item Exponent: $E = -3$, biased $= -3 + 16 = 13 = 01101_2$.
  \item Mantissa: $0100000000$ (10 bits, hidden 1 dropped).
\end{itemize}

Bit pattern:
\[
  \boxed{\texttt{1\;01101\;0100000000}}
\]

In hex:
\[
  1\,01101\,01\,00000000 = 10110101\,00000000 = \boxed{\texttt{0xB500}}
\]

\medskip\noindent\textbf{Comparison with IEEE 754 single:}
\begin{itemize}
  \item \textbf{Range:} Excess-16 gives exponents $-16$ to $+15$ (excluding
        special values), so the range is roughly $2^{-16}$ to $2^{+15} \approx
        6.1 \times 10^{-5}$ to $3.3 \times 10^{4}$. IEEE 754 single has range
        $\approx 1.2 \times 10^{-38}$ to $3.4 \times 10^{+38}$.
        The half precision range is \emph{much} smaller.
  \item \textbf{Accuracy:} 10-bit mantissa + hidden 1 = 11 bits of precision
        ($\approx 3.3$ decimal digits). IEEE 754 single has 24 bits ($\approx
        7.2$ decimal digits). Half precision is much less accurate.
\end{itemize}

%-----------------------------------------------------------------------------
\subsection*{3.28: HP 2114 Floating Point Format}
\addcontentsline{toc}{subsection}{3.28: HP 2114 Floating Point Format}
\hypertarget{sol-3.28}{}
\markright{3.28}

Format (32 bits total):
\begin{itemize}
  \item Bits 31--16: 16-bit two's complement fraction.
  \item Bits 15--8: 8-bit extension of fraction (total 24-bit fraction).
  \item Bits 7--0: 8-bit exponent in sign-magnitude with sign bit on the far
        right (bit 0).
\end{itemize}
No hidden 1.

Normalize so that $0.5 \le |f| < 1$:
\[
  -0.15625 = -0.625 \times 2^{-2}
\]

\noindent\textbf{Fraction:} $-0.625$ in 24-bit two's complement with the binary point
after the sign bit (same calculation as Problem 3.26):
\[
  -0.625 = \texttt{1011\,0000\,0000\,0000\,0000\,0000}_2 \quad\text{(24 bits)}
\]

First 16 bits (bits 31--16): $\texttt{1011\,0000\,0000\,0000}$.
Next 8 bits (bits 15--8): $\texttt{0000\,0000}$.

\noindent\textbf{Exponent:} $-2$ in sign-magnitude with sign bit on the far
right (bit 0 of the exponent byte).
Magnitude $= 2 = 0000010_2$ (7 bits). Sign $= 1$ (negative).
\[
  \text{8-bit exponent field: } \texttt{0000010\,1}_2 = \texttt{00000101}_2
\]

\noindent\textbf{32-bit pattern:}
\[
  \underbrace{\texttt{1011\,0000\,0000\,0000}}_{\text{fraction hi-16}}\;
  \underbrace{\texttt{0000\,0000}}_{\text{fraction lo-8}}\;
  \underbrace{\texttt{0000\,0101}}_{\text{exponent (sign-mag, sign on right)}}
\]

Hex:
\[
  \boxed{\texttt{0xB0000005}}
\]

\medskip\noindent\textbf{Comparison with IEEE 754 single:}
\begin{itemize}
  \item \textbf{Range:} 7-bit magnitude exponent gives range $2^{-127}$ to
        $2^{+127}$, roughly matching IEEE 754 single ($2^{-126}$ to $2^{+127}$).
  \item \textbf{Accuracy:} 24-bit two's complement fraction without hidden 1
        gives 23 bits of magnitude precision. IEEE 754 single has 23 mantissa
        bits + hidden 1 = 24 bits. The HP format has slightly less precision
        due to the lack of hidden 1 (effectively $\sim$23 vs.\ 24 bits).
  \item Both formats are 32 bits, so the HP format trades the hidden-1 trick
        for a slightly different exponent representation.
\end{itemize}
%-----------------------------------------------------------------------------
\subsection*{3.29: Half Precision Addition}
\addcontentsline{toc}{subsection}{3.29: Half Precision Addition}
\hypertarget{sol-3.29}{}
\markright{3.29}

Add $A = 2.6125 \times 10^{1} = 26.125$ and $B = 4.150390625 \times 10^{-1}
= 0.4150390625$ in 16-bit half precision (Ex.\ 3.27: 1 sign, 5-bit exponent
with excess-16 bias, 10-bit mantissa, hidden 1).

\noindent\textbf{Step 1: Convert to half precision.}
\begin{itemize}
  \item $26.125 = 11010.001_2 = 1.1010001_2 \times 2^4$.
    \[
      \text{Sign} = 0,\quad
      \text{Exp} = 4 + 16 = 20 = 10100_2,\quad
      \text{Mantissa} = 1010001000
    \]
  \item $0.4150390625 = 0.0110101_2 = 1.10101_2 \times 2^{-2}$.
    \[
      \text{Sign} = 0,\quad
      \text{Exp} = -2 + 16 = 14 = 01110_2,\quad
      \text{Mantissa} = 1010100000
    \]
\end{itemize}

\noindent\textbf{Step 2: Align exponents.}

Shift $B$'s significand right by $4 - (-2) = 6$ positions to match $A$'s
exponent of $4$:
\begin{align*}
  A &= 1.1010001000 \times 2^4 \\
  B &= 0.0000011010\;1\underbrace{00}_{\text{guard, round}} \times 2^4
       \quad(\text{sticky} = 0)
\end{align*}

\noindent\textbf{Step 3: Add significands.}
\[
  1.1010001000 + 0.0000011010 = 1.1010100010
\]
Guard $= 1$, round $= 0$, sticky $= 0$.

\noindent\textbf{Step 4: Normalize.}

Already normalized (leading 1 before the binary point).

\noindent\textbf{Step 5: Round.}

GRS $= 100$, which is exactly halfway. The LSB of the retained mantissa is $0$.
Round-to-nearest-even keeps LSB $= 0$, so we do \emph{not} round up.

Result mantissa: $1010100010$ (10 bits after hidden 1).

\noindent\textbf{Result:}
\[
  \text{Sign} = 0,\quad
  \text{Exp} = 10100_2,\quad
  \text{Mantissa} = 1010100010
\]

Bit pattern:
\[
  \texttt{0\,10100\,1010100010}
  = \texttt{01010010\,10100010}
  = \texttt{0x52A2}
\]

Decimal value:
\[
  1.1010100010_2 \times 2^4
  = 11010.100010_2
  = 26.53125
\]

The exact result is $26.125 + 0.4150390625 = 26.5400390625$. The half
precision result $26.53125$ has a small rounding error due to limited precision.

%-----------------------------------------------------------------------------
\subsection*{3.30: Half Precision Multiplication}
\addcontentsline{toc}{subsection}{3.30: Half Precision Multiplication}
\hypertarget{sol-3.30}{}
\markright{3.30}

Multiply $A = -8.0546875$ and $B = -0.179931640625$ in half precision.

\noindent\textbf{Step 1: Convert to half precision.}

\medskip\noindent\emph{Convert $A = 8.0546875$:}
\[
  8.0546875 = 1000.0000111_2 = 1.0000000111_2 \times 2^3
\]
\[
  \text{Sign} = 1,\quad
  \text{Exp} = 3 + 16 = 19 = 10011_2,\quad
  \text{Mantissa} = 0000000111
\]

\medskip\noindent\emph{Convert $B = 0.179931640625$} using repeated doubling:

\begin{center}
\begin{tabular}{rll}
\toprule
\textbf{Step} & \textbf{Computation} & \textbf{Bit} \\
\midrule
1  & $0.179931640625 \times 2 = 0.35986328125$  & 0 \\
2  & $0.35986328125  \times 2 = 0.7197265625$   & 0 \\
3  & $0.7197265625   \times 2 = 1.439453125$    & 1 \\
4  & $0.439453125    \times 2 = 0.87890625$     & 0 \\
5  & $0.87890625     \times 2 = 1.7578125$      & 1 \\
6  & $0.7578125      \times 2 = 1.515625$       & 1 \\
7  & $0.515625       \times 2 = 1.03125$        & 1 \\
8  & $0.03125        \times 2 = 0.0625$         & 0 \\
9  & $0.0625         \times 2 = 0.125$          & 0 \\
10 & $0.125          \times 2 = 0.25$           & 0 \\
11 & $0.25           \times 2 = 0.5$            & 0 \\
12 & $0.5            \times 2 = 1.0$            & 1 \\
\bottomrule
\end{tabular}
\end{center}

\[
  0.179931640625 = 0.001011100001_2 = 1.011100001_2 \times 2^{-3}
\]
\[
  \text{Sign} = 1,\quad
  \text{Exp} = -3 + 16 = 13 = 01101_2,\quad
  \text{Mantissa} = 0111000010
\]

\noindent\textbf{Step 2: Determine result sign.}
\[
  \text{Negative} \times \text{Negative} = \text{Positive} \implies \text{Sign} = 0
\]

\noindent\textbf{Step 3: Add exponents.}
\[
  E_{\text{biased}} = E_A + E_B - \text{bias} = 19 + 13 - 16 = 16 = 10000_2
\]
Equivalently, unbiased: $3 + (-3) = 0$, so biased result $= 0 + 16 = 16$.

\noindent\textbf{Step 4: Multiply significands.}
\[
  1.0000000111 \times 1.0111000010
\]

Computing the exact decimal product:
\[
  8.0546875 \times 0.179931640625 = 1.449462890625
\]

Converting to binary:
\[
  1.449462890625 \approx 1.0111001011\ldots_2 \times 2^0
\]

\noindent\textbf{Step 5: Normalize and round.}

The product is already normalized (starts with $1.\ldots$). Truncate and round
to 10-bit mantissa:
\[
  1.0111001011_2 \times 2^0
\]

\noindent\textbf{Step 6: Check overflow/underflow.}

Biased exponent $= 16$, which is in range $[1, 30]$. No overflow or underflow.

\noindent\textbf{Result:}
\[
  \text{Sign} = 0,\quad
  \text{Exp} = 10000_2,\quad
  \text{Mantissa} = 0111001011
\]

Bit pattern:
\[
  \texttt{0\,10000\,0111001011} = \texttt{0x41CB}
\]

Decimal value:
\[
  1.0111001011_2 \times 2^0 \approx 1.44921875
\]

Exact product: $8.0546875 \times 0.179931640625 = 1.449462890625$.
Half precision result: $1.44921875$. Relative error $\approx 0.017\%$.

%-----------------------------------------------------------------------------
\subsection*{3.31: Half Precision Division}
\addcontentsline{toc}{subsection}{3.31: Half Precision Division}
\hypertarget{sol-3.31}{}
\markright{3.31}

Divide $A = 86.25$ by $B = -4.875$ in half precision.

\noindent\textbf{Step 1: Convert to half precision.}
\begin{itemize}
  \item $86.25 = 1010110.01_2 = 1.01011001_2 \times 2^6$.
    \[
      \text{Sign} = 0,\quad
      \text{Exp} = 6 + 16 = 22 = 10110_2,\quad
      \text{Mantissa} = 0101100100
    \]
  \item $4.875 = 100.111_2 = 1.00111_2 \times 2^2$.
    \[
      \text{Sign} = 1,\quad
      \text{Exp} = 2 + 16 = 18 = 10010_2,\quad
      \text{Mantissa} = 0011100000
    \]
\end{itemize}

\noindent\textbf{Step 2: Determine result sign.}
\[
  \text{Positive} \div \text{Negative} = \text{Negative} \implies \text{Sign} = 1
\]

\noindent\textbf{Step 3: Subtract exponents.}
\[
  E_{\text{biased}} = E_A - E_B + \text{bias} = 22 - 18 + 16 = 20 = 10100_2
\]
Equivalently, unbiased: $6 - 2 = 4$, so biased result $= 4 + 16 = 20$.

\noindent\textbf{Step 4: Divide significands.}

The exact decimal quotient is:
\[
  86.25 \div 4.875 = 17.6923076923\ldots
\]

Converting to binary:
\[
  17.6923\ldots \approx 10001.1011_2 = 1.00011011_2 \times 2^4
\]
The exponent from Step~3 is already $4$, so the significand quotient is
$\approx 1.0001101100_2$.

\noindent\textbf{Step 5: Normalize and round.}

Already normalized. Rounding the 11-bit significand to 10-bit mantissa
with round-to-nearest-even gives:
\[
  1.0001101100_2 \times 2^4
\]

\noindent\textbf{Result:}
\[
  \text{Sign} = 1,\quad
  \text{Exp} = 10100_2 = 20,\quad
  \text{Mantissa} = 0001101100
\]

Bit pattern:
\[
  \texttt{1\,10100\,0001101100}
  = \texttt{1101\,0000\,0110\,1100}
  = \texttt{0xD06C}
\]

Decimal value:
\[
  -1.0001101100_2 \times 2^4
  = -10001.101100_2
  = -(16 + 1 + 0.5 + 0.125 + 0.0625)
  = -17.6875
\]

Exact quotient: $86.25 / 4.875 = 17.6923\ldots$, represented as $-17.6875$.
Relative error: $|17.6923 - 17.6875| / 17.6923 \approx 0.027\%$.

%-----------------------------------------------------------------------------
\subsection*{3.32: FP Addition Associativity (Left)}
\addcontentsline{toc}{subsection}{3.32: FP Addition Associativity (Left)}
\hypertarget{sol-3.32}{}
\markright{3.32}

Compute $(3.984375 \times 10^{-1} + 3.4375 \times 10^{-1}) + 1.771 \times 10^{3}$
in half precision.

\noindent\textbf{First addition:} $0.3984375 + 0.34375$.

Convert both values:
\begin{align*}
  0.3984375 &= 0.0110011_2 = 1.10011_2 \times 2^{-2} \\
  &\quad \text{Sign} = 0,\; \text{Exp} = -2 + 16 = 14 = 01110_2,\;
    \text{Mantissa} = 1001100000 \\[4pt]
  0.34375 &= 0.01011_2 = 1.011_2 \times 2^{-2} \\
  &\quad \text{Sign} = 0,\; \text{Exp} = 14 = 01110_2,\;
    \text{Mantissa} = 0110000000
\end{align*}

Exponents match, so add directly:
\[
  1.1001100000 + 1.0110000000 = 10.1111100000
\]

Normalize:
\[
  10.1111100000 \times 2^{-2}
  = 1.01111100000 \times 2^{-1}
\]
Exponent becomes $-1 + 16 = 15$. Mantissa $= 0111110000$ (fits in 10 bits).
GRS $= 000$, no rounding needed.

Intermediate result:
\[
  1.0111110000 \times 2^{-1} = 0.10111110000_2 = 0.7421875
\]

\noindent\textbf{Second addition:} $0.7421875 + 1771$.

Convert $1771$:
\[
  1771 = 11011101011_2 = 1.1011101011_2 \times 2^{10}
\]
\[
  \text{Sign} = 0,\quad
  \text{Exp} = 10 + 16 = 26 = 11010_2,\quad
  \text{Mantissa} = 1011101011
\]

Align: shift the smaller number right by $10 - (-1) = 11$ positions:
\[
  1.0111110000 \times 2^{-1} \to 0.00000000001\;01111\ldots \times 2^{10}
\]

After shifting 11 places, the significand is essentially 0 at 10-bit precision
(all shifted bits become guard/round/sticky). Sticky bit $= 1$ from shifted-out
bits.

Sum:
\[
  \approx 1.1011101011 \times 2^{10}
\]

GRS $= 001$: round down since $< 0.5$ ULP.

\noindent\textbf{Result:}
\[
  1.1011101011 \times 2^{10},\quad
  \text{Sign} = 0,\quad
  \text{Exp} = 11010_2,\quad
  \text{Mantissa} = 1011101011
\]

Bit pattern: $\texttt{0\,11010\,1011101011}$.

Decimal: $\mathbf{1771.0}$ (the small value was lost).

%-----------------------------------------------------------------------------
\subsection*{3.33: FP Addition Associativity (Right)}
\addcontentsline{toc}{subsection}{3.33: FP Addition Associativity (Right)}
\hypertarget{sol-3.33}{}
\markright{3.33}

Compute $3.984375 \times 10^{-1} + (3.4375 \times 10^{-1} + 1.771 \times 10^{3})$
in half precision.

\noindent\textbf{First addition:} $0.34375 + 1771$.

\[
  0.34375 = 1.011_2 \times 2^{-2},\qquad
  1771 = 1.1011101011_2 \times 2^{10}
\]

Align: shift $0.34375$ right by $10 - (-2) = 12$ positions to match exponent
$2^{10}$. The entire significand is shifted out, so the value is lost.

Result:
\[
  \approx 1.1011101011 \times 2^{10} = 1771.0
\]

\noindent\textbf{Second addition:} $0.3984375 + 1771$.

\[
  0.3984375 = 1.10011_2 \times 2^{-2}
\]

Align: shift right by $10 - (-2) = 12$ positions. Again, the entire significand
is shifted out and the value is lost.

Result:
\[
  \approx 1.1011101011 \times 2^{10} = 1771.0
\]

\noindent\textbf{Final result:} $\mathbf{1771.0}$, same as Problem~3.32.

%-----------------------------------------------------------------------------
\subsection*{3.34: FP Addition Associativity Check}
\addcontentsline{toc}{subsection}{3.34: FP Addition Associativity Check}
\hypertarget{sol-3.34}{}
\markright{3.34}

\textbf{Yes}, in this case the results of 3.32 and 3.33 agree: both give
$1771.0$.

However, this is somewhat coincidental. Both small values are completely
swamped by the large value due to the limited 11-bit significand. In general,
floating-point addition is \textbf{not associative}. A different set of values
could produce different results depending on grouping, since rounding occurs at
each step.

If computed exactly:
\[
  0.3984375 + 0.34375 + 1771 = 1771.7421875
\]
This differs from $1771.0$. The loss of the small addends demonstrates the
precision limitations of half-precision floating point.
%-----------------------------------------------------------------------------
\subsection*{3.35: FP Multiply Associativity (Left)}
\addcontentsline{toc}{subsection}{3.35: FP Multiply Associativity (Left)}
\hypertarget{sol-3.35}{}
\markright{3.35}

Compute $(3.41796875 \times 10^{-3} \times 6.34765625 \times 10^{-3})
\times 1.05625 \times 10^{2}$ in half precision (1 sign, 5-bit exponent with
excess-16 bias, 10-bit mantissa, hidden~1).

\medskip\noindent\textbf{Step 1: Convert each value to half precision.}

\begin{itemize}
  \item $A = 0.00341796875$.
        Multiply by $2^{9}$: $0.00341796875 \times 512 = 1.75 = 1.11_2$.
        So $A = 1.1100000000_2 \times 2^{-9}$.
        Sign $= 0$, biased exponent $= -9 + 16 = 7 = 00111_2$,
        mantissa $= 1100000000$.
  \item $B = 0.00634765625$.
        Multiply by $2^{8}$: $0.00634765625 \times 256 = 1.625 = 1.101_2$.
        So $B = 1.1010000000_2 \times 2^{-8}$.
        Sign $= 0$, biased exponent $= -8 + 16 = 8 = 01000_2$,
        mantissa $= 1010000000$.
  \item $C = 105.625 = 1101001.101_2 = 1.1010011010_2 \times 2^{6}$.
        Sign $= 0$, biased exponent $= 6 + 16 = 22 = 10110_2$,
        mantissa $= 1010011010$.
\end{itemize}

\noindent\textbf{Step 2: First multiplication $A \times B$.}

Result sign: positive.  Exponent sum: $(-9) + (-8) = -17$.

Multiply significands:
\begin{verbatim}
        1.1100
      x 1.1010
      --------
        1.1100      (1.11 x 1)
    +   0.1110 0    (1.11 x 0.1)
    +   0.0011 10   (1.11 x 0.001)
    ----------------
       10.1101 10
\end{verbatim}

\noindent So the product is $10.110110_2 \times 2^{-17}$.  Normalize:
\[
  1.0110110000_2 \times 2^{-16}.
\]
Biased exponent $= -16 + 16 = 0$.  Since biased~0 is reserved for
denormals/zero, this is an \textbf{underflow}.

Represent as a denormalized number (biased exponent~$= 0$, implied leading~0,
effective exponent $= 1 - 16 = -15$):
\[
  1.011011_2 \times 2^{-16} = 0.1011011000_2 \times 2^{-15}.
\]
Denormal mantissa bits (no hidden~1): $1011011000$.  GRS $= 000$ (exact fit),
so no rounding is needed.

Stored value: $\texttt{0\;00000\;1011011000}$.
\[
  A \times B = 0.1011011_2 \times 2^{-15}
  = \frac{91}{128} \times 2^{-15}
  \approx 2.170 \times 10^{-5}.
\]

\noindent\textbf{Step 3: Second multiplication $(A \times B) \times C$.}

The first product is denormalized with significand $0.1011011000$ and effective
exponent $-15$.  Multiply by $C = 1.1010011010_2 \times 2^{6}$:
\[
  \text{Exponent: } {-15} + 6 = -9.
\]
Significand product (using exact arithmetic on the stored bits):
\[
  \frac{91}{128} \times \frac{845}{512}
  = \frac{76895}{65536}
  = 1.0010110001011111_2.
\]
This is already normalized ($1 \le 1.0010\ldots < 2$), so the exponent remains
$-9$.  Biased exponent $= -9 + 16 = 7 = 00111_2$.

Extract the 10-bit mantissa with GRS:
\[
  1.\underbrace{0010110001}_{10\text{ bits}}\underbrace{0\;1\;1}_{\text{G R S}}111\ldots
\]
GRS $= 011$ (less than halfway) $\Rightarrow$ truncate.
Mantissa $= 0010110001$.

\medskip\noindent\textbf{Result:}

\[
  \boxed{\texttt{0\;00111\;0010110001}}
\]

Decimal value:
\[
  \bigl(1 + 177/1024\bigr) \times 2^{-9}
  = \frac{1201}{1024} \times \frac{1}{512}
  \approx \mathbf{0.002291} .
\]

Exact answer: $0.00341796875 \times 0.00634765625 \times 105.625
= 0.0022916\ldots$

%-----------------------------------------------------------------------------
\subsection*{3.36: FP Multiply Associativity (Right)}
\addcontentsline{toc}{subsection}{3.36: FP Multiply Associativity (Right)}
\hypertarget{sol-3.36}{}
\markright{3.36}

Compute $3.41796875 \times 10^{-3} \times (6.34765625 \times 10^{-3}
\times 1.05625 \times 10^{2})$ in half precision.

From 3.35, the half-precision representations are:
$A = 1.1100000000_2 \times 2^{-9}$,\;
$B = 1.1010000000_2 \times 2^{-8}$,\;
$C = 1.1010011010_2 \times 2^{6}$.

\medskip\noindent\textbf{Step 1: First multiplication $B \times C$.}

Result sign: positive.  Exponent sum: $(-8) + 6 = -2$.

Multiply significands:
\[
  \frac{13}{8} \times \frac{845}{512}
  = \frac{10985}{4096}
  = 2.681884765625.
\]
In binary: $10985_{10} = 10101011101001_2$, so
\[
  \frac{10985}{4096} = 10.101011101001_2.
\]
Normalize: $1.0101011101001_2 \times 2^{-1}$.
Biased exponent $= -1 + 16 = 15 = 01111_2$.

Extract 10-bit mantissa with GRS:
\[
  1.\underbrace{0101011101}_{10}\underbrace{0\;0\;1}_{\text{G R S}}.
\]
GRS $= 001$ (less than halfway) $\Rightarrow$ truncate.
Mantissa $= 0101011101$.

Stored $B \times C$: $\texttt{0\;01111\;0101011101}$.
\[
  B \times C = 1.0101011101_2 \times 2^{-1}
  = \frac{1373}{2048}
  = 0.67041015625.
\]

\noindent\textbf{Step 2: Second multiplication $A \times (B \times C)$.}

Exponent sum: $(-9) + (-1) = -10$.

Multiply significands:
\[
  \frac{7}{4} \times \frac{1373}{1024}
  = \frac{9611}{4096}
  = 2.346435546875.
\]
In binary: $9611_{10} = 10010110001011_2$, so
\[
  \frac{9611}{4096} = 10.010110001011_2.
\]
Normalize: $1.0010110001011_2 \times 2^{-9}$.
Biased exponent $= -9 + 16 = 7 = 00111_2$.

Extract 10-bit mantissa with GRS:
\[
  1.\underbrace{0010110001}_{10}\underbrace{0\;1\;1}_{\text{G R S}}.
\]
GRS $= 011$ (less than halfway) $\Rightarrow$ truncate.
Mantissa $= 0010110001$.

\medskip\noindent\textbf{Result:}

\[
  \boxed{\texttt{0\;00111\;0010110001}}
\]

Decimal value:
\[
  \bigl(1 + 177/1024\bigr) \times 2^{-9}
  = \frac{1201}{524288}
  \approx \mathbf{0.002291}.
\]

%-----------------------------------------------------------------------------
\subsection*{3.37: FP Multiply Associativity Check}
\addcontentsline{toc}{subsection}{3.37: FP Multiply Associativity Check}
\hypertarget{sol-3.37}{}
\markright{3.37}

\textbf{Yes}, for these particular values the results of 3.35 and 3.36
\textbf{agree}: both produce $\texttt{0\;00111\;0010110001} \approx 0.002291$.

This happens because the intermediate product $A \times B$ in 3.35, despite
underflowing to a denormalized number, is \emph{exactly} representable (GRS
$= 000$, no rounding). Therefore no precision is lost in the denormal step, and
both orderings produce the same final rounded result.

In general, however, floating-point multiplication is \textbf{not associative}.
If the denormal representation had required rounding, or if different
intermediate results triggered different rounding decisions, the two orderings
would produce different answers.

The exact result is $0.00341796875 \times 0.00634765625 \times 105.625
= 0.0022916\ldots$, while both half-precision computations give
$\approx 0.002291$.

%-----------------------------------------------------------------------------
\subsection*{3.38: FP Distributivity (Factored Form)}
\addcontentsline{toc}{subsection}{3.38: FP Distributivity (Factored Form)}
\hypertarget{sol-3.38}{}
\markright{3.38}

Compute $1.666015625 \times (19760 + (-19744))$ in half precision.

\medskip\noindent\textbf{Step 1: Convert to half precision.}

\begin{itemize}
  \item $19760 = 100110100110000_2 = 1.00110100110000_2 \times 2^{14}$.
        Biased exponent $= 14 + 16 = 30 = 11110_2$.
        10-bit mantissa: $0011010011$ (remaining bits $0000$; GRS $= 000$, exact).
  \item $19744 = 100110100100000_2 = 1.00110100100000_2 \times 2^{14}$.
        Biased exponent $= 30 = 11110_2$.
        10-bit mantissa: $0011010010$ (exact fit).
  \item $1.666015625 = 1.1010101010_2 \times 2^{0}$.
        Biased exponent $= 0 + 16 = 16 = 10000_2$.
        Mantissa: $1010101010$ (exact).
\end{itemize}

\noindent\textbf{Step 2: Subtract inside the parentheses.}

Both operands share exponent $2^{14}$.  Subtract significands directly:
\[
  1.0011010011_2 - 1.0011010010_2 = 0.0000000001_2 \times 2^{14}.
\]
Normalize: $1.0000000000_2 \times 2^{4}$.
Biased exponent $= 4 + 16 = 20 = 10100_2$, mantissa $= 0000000000$.
GRS $= 000$; no rounding.  The subtraction is \textbf{exact}: $19760 - 19744 = 16$.

\medskip\noindent\textbf{Step 3: Multiply $1.666015625 \times 16$.}

\[
  1.1010101010_2 \times 2^{0} \;\times\; 1.0000000000_2 \times 2^{4}.
\]
Exponent sum: $0 + 4 = 4$.  Significand product:
\[
  1.1010101010 \times 1.0 = 1.1010101010.
\]
Already normalized.  Biased exponent $= 4 + 16 = 20 = 10100_2$.
Mantissa $= 1010101010$.  GRS $= 000$; no rounding needed.

\medskip\noindent\textbf{Result:}
\[
  \boxed{\texttt{0\;10100\;1010101010}}
\]
Decimal:
\[
  1.1010101010_2 \times 2^4 = 11010.101010_2 = \mathbf{26.65625}.
\]

The exact answer is $1.666015625 \times 16 = 26.65625$, so the half-precision
result is \textbf{exact}.

%-----------------------------------------------------------------------------
\subsection*{3.39: FP Distributivity (Expanded Form)}
\addcontentsline{toc}{subsection}{3.39: FP Distributivity (Expanded Form)}
\hypertarget{sol-3.39}{}
\markright{3.39}

Compute $(1.666015625 \times 19760) + (1.666015625 \times (-19744))$ in half
precision.

Using the representations from 3.38:
$P = 1.1010101010_2 \times 2^{0}$,
$Q = 1.0011010011_2 \times 2^{14}$ (for 19760),
$R = 1.0011010010_2 \times 2^{14}$ (for 19744).

\medskip\noindent\textbf{Step 1: First product $P \times Q = 1.666015625 \times
19760$.}

Exponent sum: $0 + 14 = 14$.  Significand product:
\[
  \frac{853}{512} \times \frac{1235}{1024}
  = \frac{1{,}053{,}455}{524{,}288}
  = 2.009305\ldots
\]
Normalize: $1.00000010011000011110\ldots_2 \times 2^{15}$.
Biased exponent $= 15 + 16 = 31 = 11111_2$.

10-bit mantissa: $0000000100$.
GRS from remaining bits ($1100001111\ldots$): $\text{G}=1,\;\text{R}=1,\;
\text{S}=1$.
Round up: mantissa $= 0000000101$.

If biased exponent 31 is valid (simple excess-16), the stored value is:
\[
  \texttt{0\;11111\;0000000101}, \quad
  \text{value} = \bigl(1 + \tfrac{5}{1024}\bigr) \times 2^{15} = 32{,}928.
\]
(The exact product is $32{,}920.46875$, so rounding error $\approx 7.5$.)

\medskip\noindent\textbf{Step 2: Second product $P \times R = 1.666015625 \times
19744$.}

Exponent sum: $0 + 14 = 14$.  Significand product:
\[
  \frac{853}{512} \times \frac{617}{512}
  = \frac{526{,}301}{262{,}144}
  = 2.007678\ldots
\]
Normalize: $1.00000000111110111010\ldots_2 \times 2^{15}$.
Biased exponent $= 31$.

10-bit mantissa: $0000000011$.
GRS $= 111$.
Round up: mantissa $= 0000000100$.

Stored value:
\[
  \texttt{1\;11111\;0000000100}, \quad
  \text{value} = -\bigl(1 + \tfrac{4}{1024}\bigr) \times 2^{15} = -32{,}896.
\]
(The exact product is $-32{,}893.8125$.)

\medskip\noindent\textbf{Step 3: Add the two products.}

Both magnitudes share exponent $2^{15}$.  Subtract significands:
\[
  1.0000000101 - 1.0000000100 = 0.0000000001 \times 2^{15}.
\]
Normalize: $1.0 \times 2^{5}$.
Biased exponent $= 5 + 16 = 21 = 10101_2$, mantissa $= 0000000000$.

\medskip\noindent\textbf{Result:}
\[
  \boxed{\texttt{0\;10101\;0000000000}}
\]
Decimal:
\[
  1.0 \times 2^5 = \mathbf{32.0}.
\]
The exact answer is $26.65625$, so the half-precision result has large error
($\approx 20\%$).

\medskip\noindent\textbf{Why the error is so large.}
The two intermediate products ($\approx 32{,}920$ and $\approx 32{,}894$)
are close in magnitude, so their difference is small.  At exponent $2^{15}$,
the spacing between consecutive half-precision values is
$2^{15}/2^{10} = 32$, which is \emph{larger} than the true difference of
$26.65625$.  The individual rounding errors ($+7.53$ and $-2.19$) do not cancel;
instead they dominate the result.  This is \textbf{catastrophic cancellation}.

\medskip\noindent\textit{Note:} If biased exponent 31 is reserved for
$\infty$/NaN (following IEEE 754 conventions), both products overflow to
$\pm\infty$ and the sum is $+\infty + (-\infty) = \text{NaN}$---an even worse
outcome.

%-----------------------------------------------------------------------------
\subsection*{3.40: FP Distributivity Check}
\addcontentsline{toc}{subsection}{3.40: FP Distributivity Check}
\hypertarget{sol-3.40}{}
\markright{3.40}

\textbf{No}, the results of 3.38 and 3.39 do \textbf{not} agree.

\begin{itemize}
  \item 3.38 (factored form): $1.666015625 \times 16 = 26.65625$ (\textbf{exact}).
        Bit pattern: $\texttt{0\;10100\;1010101010}$.
  \item 3.39 (expanded form): $32{,}928 + (-32{,}896) = 32.0$
        (\textbf{$\approx 20\%$ error}).
        Bit pattern: $\texttt{0\;10101\;0000000000}$.
\end{itemize}

This is a classic demonstration that the \textbf{distributive law does not hold}
for floating-point arithmetic.  The expanded form computes two large
intermediate products ($\approx 32{,}900$) whose difference is small
($\approx 27$).  Because the 10-bit mantissa cannot represent both the large
magnitude and the fine detail simultaneously, the rounding errors in the two
products are comparable to the true difference.  Subtracting the rounded values
produces \textbf{catastrophic cancellation}.

The factored form avoids this by computing the small difference ($19760 -
19744 = 16$) first---an exact operation---and then multiplying, preserving full
precision throughout.

%-----------------------------------------------------------------------------
\subsection*{3.41: IEEE 754 Representation of $-1/4$}
\addcontentsline{toc}{subsection}{3.41: IEEE 754 Representation of $-1/4$}
\hypertarget{sol-3.41}{}
\markright{3.41}

\[
  -\frac{1}{4} = -0.25 = -0.01_2 = -1.0_2 \times 2^{-2}
\]

IEEE 754 single precision:
\begin{itemize}
  \item Sign: $1$.
  \item Exponent:
    \[
      -2 + 127 = 125 = 01111101_2
    \]
  \item Mantissa: $00000000000000000000000$ (all zeros; the hidden 1 provides
        the $1.0$).
\end{itemize}

Bit pattern:
\[
  \boxed{\texttt{1\;01111101\;00000000000000000000000}}
\]
Hex: $\texttt{0xBE800000}$.

\textbf{Yes}, $-1/4$ can be represented \textbf{exactly} in IEEE 754. It is
a negative power of 2, so it has an exact binary representation with a finite
number of bits.

%-----------------------------------------------------------------------------
\subsection*{3.42: Repeated Addition vs Multiplication}
\addcontentsline{toc}{subsection}{3.42: Repeated Addition vs Multiplication}
\hypertarget{sol-3.42}{}
\markright{3.42}

\noindent\textbf{Adding $-1/4$ to itself 4 times:}
\begin{align*}
  &(-1/4) + (-1/4) = -1/2 \quad\text{(exact)} \\
  &(-1/2) + (-1/4) = -3/4 \quad\text{(exact)} \\
  &(-3/4) + (-1/4) = -1   \quad\text{(exact)}
\end{align*}

All intermediate results are exact in IEEE 754 (they are all dyadic rationals
with small denominators).

\noindent\textbf{$-1/4 \times 4$:}
\[
  -\frac{1}{4} \times 4 = -1 \quad\text{(exact, since 4 is a power of 2)}
\]

\textbf{Yes}, they are the same: both give $\mathbf{-1}$, which is also the
mathematically correct answer.

In this case, all values involved are exactly representable and all operations
produce exact results. This is a special case --- for most values, repeated
addition and multiplication would not agree due to accumulated rounding errors.

%-----------------------------------------------------------------------------
\subsection*{3.43: Binary Fraction of $1/3$}
\addcontentsline{toc}{subsection}{3.43: Binary Fraction of $1/3$}
\hypertarget{sol-3.43}{}
\markright{3.43}

$1/3$ in binary (24 bits):
\[
  \frac{1}{3} = 0.\overline{01}_2 = 0.010101010101010101010101\ldots_2
\]

The pattern $01$ repeats forever. With 24 bits:
\[
  \texttt{010101010101010101010101}_2
\]

This representation is \textbf{not exact}. The value $1/3$ is a repeating
fraction in binary (just as $1/3 = 0.\overline{3}$ repeats in decimal). No
finite number of binary digits can represent it exactly.

The 24-bit approximation gives:
\[
  \sum_{k=1}^{12} 2^{-2k}
  = \frac{1}{4}\cdot\frac{1-(1/4)^{12}}{1-1/4}
  = \frac{1}{3}\bigl(1 - 4^{-12}\bigr)
  \approx 0.33333325
\]

%-----------------------------------------------------------------------------
\subsection*{3.44: BCD Fraction of $1/3$}
\addcontentsline{toc}{subsection}{3.44: BCD Fraction of $1/3$}
\hypertarget{sol-3.44}{}
\markright{3.44}

In BCD (Binary Coded Decimal), each decimal digit uses 4 bits. With 24 bits,
we can store $24/4 = 6$ decimal digits.
\[
  \frac{1}{3} = 0.\overline{3}_{10} = 0.333333\ldots_{10}
\]

With 6 BCD digits: $\texttt{0011\,0011\,0011\,0011\,0011\,0011}$ (each
$\texttt{0011}$ encodes the digit 3).

This represents $0.333333$, which is \textbf{not exact}. The value $1/3$ is a
repeating fraction in base 10, so no finite number of BCD digits can represent
it exactly.

%-----------------------------------------------------------------------------
\subsection*{3.45: Base-16 Fraction of $1/3$}
\addcontentsline{toc}{subsection}{3.45: Base-16 Fraction of $1/3$}
\hypertarget{sol-3.45}{}
\markright{3.45}

In base 16, each digit uses 4 bits. With 24 bits, we have $24/4 = 6$
hexadecimal digit positions.

$1/3$ in base 16: $1/3 = 0.\overline{5}_{16}$. Verify:
\[
  5 \times \bigl(16^{-1} + 16^{-2} + 16^{-3} + \cdots\bigr)
  = 5 \times \frac{1/16}{1 - 1/16}
  = \frac{5}{15}
  = \frac{1}{3} \quad\checkmark
\]

With 6 hex digits (24 bits): $\texttt{0101\,0101\,0101\,0101\,0101\,0101}$.

This represents $0.555555_{16}$:
\[
  5\bigl(16^{-1} + 16^{-2} + \cdots + 16^{-6}\bigr)
  = 5 \times \frac{1/16 - 16^{-7}}{15/16}
  = \frac{1}{3}\bigl(1 - 16^{-6}\bigr)
\]

This is \textbf{not exact}. The digit 5 repeats infinitely in base 16, just
as it does in other bases that are not multiples of 3.

%-----------------------------------------------------------------------------
\subsection*{3.46: Base-30 Fraction of $1/3$}
\addcontentsline{toc}{subsection}{3.46: Base-30 Fraction of $1/3$}
\hypertarget{sol-3.46}{}
\markright{3.46}

In base 30, each digit requires $\lceil \log_2 30 \rceil = 5$ bits. With 20
bits, we have $20/5 = 4$ base-30 digit positions.

$1/3$ in base 30:
\[
  \frac{1}{3} = \frac{10}{30} = 0.\text{A}_{30}
\]
where A represents the digit 10. This is \textbf{exact} in a single base-30
digit, since $10/30 = 1/3$ exactly.

With 20 bits (4 digit positions): $\texttt{01010\,00000\,00000\,00000}$
(digit 10 = $01010_2$ followed by three zero digits).

This representation \textbf{is exact}. This works because 30 is divisible by 3,
so $1/3$ has a terminating representation in base 30.

In general, a fraction $p/q$ has a terminating representation in base $b$ if and
only if all prime factors of $q$ are also factors of $b$. Since $3 \mid 30$, the
fraction $1/3$ terminates in base 30.

%-----------------------------------------------------------------------------
\subsection*{3.47: SIMD FIR Filter Implementation}
\addcontentsline{toc}{subsection}{3.47: SIMD FIR Filter Implementation}
\hypertarget{sol-3.47}{}
\markright{3.47}

The 4-tap FIR filter computes:
\begin{verbatim}
    sig_out[i] = sig_in[i-3]*f[0] + sig_in[i-2]*f[1]
               + sig_in[i-1]*f[2] + sig_in[i]*f[3]
\end{verbatim}
with coefficients $f = \{-1, 2, -1, 2\}$ and 16-bit values.

\medskip\noindent\textbf{SIMD strategy with 128-bit registers:}

Each 128-bit register can hold $128/16 = 8$ values of 16-bit data.

\begin{enumerate}
  \item \textbf{Load coefficient vectors.} Create four 128-bit registers, each
        containing 8 copies of one coefficient:
        \begin{itemize}
          \item \texttt{vf0} $= \{-1, -1, -1, -1, -1, -1, -1, -1\}$
          \item \texttt{vf1} $= \{2, 2, 2, 2, 2, 2, 2, 2\}$
          \item \texttt{vf2} $= \{-1, -1, -1, -1, -1, -1, -1, -1\}$
          \item \texttt{vf3} $= \{2, 2, 2, 2, 2, 2, 2, 2\}$
        \end{itemize}

  \item \textbf{Process 8 outputs at a time.} For each block of 8 output
        positions $i, i+1, \ldots, i+7$:
        \begin{itemize}
          \item Load \texttt{v0} $= \{\text{sig\_in}[i-3], \ldots,
                \text{sig\_in}[i+4]\}$ (8 values starting at $i-3$).
          \item Load \texttt{v1} $= \{\text{sig\_in}[i-2], \ldots,
                \text{sig\_in}[i+5]\}$ (8 values starting at $i-2$).
          \item Load \texttt{v2} $= \{\text{sig\_in}[i-1], \ldots,
                \text{sig\_in}[i+6]\}$ (8 values starting at $i-1$).
          \item Load \texttt{v3} $= \{\text{sig\_in}[i], \ldots,
                \text{sig\_in}[i+7]\}$ (8 values starting at $i$).
        \end{itemize}

  \item \textbf{SIMD multiply-accumulate.} Use packed 16-bit multiply and add
        instructions:
\begin{verbatim}
    vr  = pmul(v0, vf0)       # 8 parallel multiplies
    vr  = pmadd(v1, vf1, vr)  # 8 parallel multiply-adds
    vr  = pmadd(v2, vf2, vr)  # 8 parallel multiply-adds
    vr  = pmadd(v3, vf3, vr)  # 8 parallel multiply-adds
\end{verbatim}

  \item \textbf{Store result.} Write the 8 results to
        \texttt{sig\_out[i..i+7]} with a single 128-bit store.

  \item \textbf{Advance.} Increment $i$ by 8 and repeat until all 125 outputs
        ($i = 3$ to $127$) are computed.
\end{enumerate}

\medskip\noindent\textbf{Assumptions:}
\begin{itemize}
  \item The processor supports packed 16-bit integer multiply (\texttt{pmul})
        and multiply-accumulate (\texttt{pmadd}) operating on 128-bit registers.
  \item Unaligned loads are supported (or alignment is handled), since \texttt{v0}
        through \texttt{v3} overlap with offsets of 2 bytes (one 16-bit element).
  \item Overflow is handled by saturating arithmetic or by using 32-bit
        accumulators (some SIMD ISAs widen the multiply result).
\end{itemize}

\medskip\noindent\textbf{Benefits:}
\begin{itemize}
  \item 8$\times$ parallelism per SIMD instruction reduces the loop iteration
        count from 125 to $\lceil 125/8 \rceil = 16$.
  \item The number of memory operations is minimized: 4 loads + 1 store per
        8 outputs, and consecutive iterations share most of the loaded data (the
        windows overlap by 7 elements, so data can be reused via register
        shifting).
  \item For further optimization, a sliding window approach can reuse previously
        loaded data by shifting registers rather than reloading, reducing memory
        bandwidth.
\end{itemize}
